\documentclass[a4paper,12pt]{article}

% ------------------------------------------------
% Кодировка, язык, математика
% ------------------------------------------------
\usepackage[T2A]{fontenc}
\usepackage[utf8]{inputenc}
\usepackage[russian]{babel}
\usepackage{amsmath,amssymb,mathtools}
\usepackage{icomma}

% ------------------------------------------------
% Шрифт и типографика
% ------------------------------------------------
\usepackage{helvet}
\renewcommand{\familydefault}{\sfdefault}
\usepackage{microtype}
\usepackage{setspace}
\onehalfspacing
\usepackage{parskip}
\setlength{\parindent}{0pt}

% ------------------------------------------------
% Страница
% ------------------------------------------------
\usepackage[
  a4paper,
  left=20mm,
  right=20mm,
  top=20mm,
  bottom=22mm,
  headheight=15pt
]{geometry}

% ------------------------------------------------
% Графика и таблицы
% ------------------------------------------------
\usepackage{tikz}
\usetikzlibrary{calc,arrows.meta,positioning,shapes.geometric}
\usepackage{pgfplots}
\pgfplotsset{compat=1.18}
\usepackage{tabularx,booktabs}
\usepackage{enumitem}
\usepackage{needspace}
\usepackage{xcolor}

% ------------------------------------------------
% Колонтитулы
% ------------------------------------------------
\usepackage{fancyhdr}

% ------------------------------------------------
% Палитра
% ------------------------------------------------
\definecolor{accent}{HTML}{008080}
\definecolor{darkaccent}{HTML}{004D4D}
\definecolor{textgray}{HTML}{333333}
\definecolor{linegray}{HTML}{A8B0B0}

\color{textgray}

\pagestyle{fancy}
\fancyhf{}
\fancyhead[L]{\small Анастасия Павлова}
\fancyhead[R]{\small Графики функций}
\fancyfoot[C]{\small\thepage}
\renewcommand{\headrulewidth}{0.3pt}
\renewcommand{\footrulewidth}{0pt}

% ------------------------------------------------
% Служебные команды
% ------------------------------------------------
\newcommand{\topic}[1]{%
  \Needspace{7\baselineskip}%
  \vspace{0.45em}%
  {\Large\bfseries\color{darkaccent} #1}\par
  \vspace{0.2em}%
  {\color{accent}\rule{\linewidth}{0.7pt}}\par
  \vspace{0.35em}%
}

\newcommand{\keyidea}[1]{%
  \begin{center}
    \begin{minipage}{0.92\linewidth}
      \centering\color{darkaccent}\bfseries #1
    \end{minipage}
  \end{center}
}

\newcommand{\checkitem}{\item[$\square$]}

% ------------------------------------------------
% Единые стили блок-схем
% ------------------------------------------------
\tikzset{
  flowstart/.style={
    ellipse,
    draw=darkaccent,
    line width=0.9pt,
    align=center,
    minimum width=4.2cm,
    minimum height=1.0cm,
    inner sep=4pt,
    font=\small
  },
  flowaction/.style={
    rectangle,
    rounded corners=3pt,
    draw=accent,
    line width=0.9pt,
    align=center,
    text width=8.8cm,
    minimum height=1.0cm,
    inner sep=5pt,
    font=\small
  },
  flowside/.style={
    rectangle,
    rounded corners=3pt,
    draw=accent,
    line width=0.9pt,
    align=center,
    text width=5.0cm,
    minimum height=1.0cm,
    inner sep=5pt,
    font=\small
  },
  flowdecision/.style={
    diamond,
    aspect=2.7,
    draw=darkaccent,
    line width=0.9pt,
    align=center,
    text width=5.0cm,
    inner sep=2pt,
    font=\small
  },
  flowarrow/.style={
    -{Latex[length=3mm,width=2mm]},
    line width=0.9pt,
    draw=textgray
  }
}

\begin{document}

% ================================================================
% ТИТУЛЬНЫЙ ЛИСТ
% ================================================================
\begin{titlepage}
\thispagestyle{empty}
\centering

\vspace*{2.4cm}

{\Large\color{darkaccent}\bfseries Чек-лист\par}
\vspace{0.7cm}
{\Huge\bfseries Графики функций\par}
\vspace{0.25cm}
{\LARGE\bfseries Узловые точки и точки пересечения\par}

\vspace{1.5cm}
{\large ЕГЭ, профильная математика\par}

\vfill

{\large\bfseries Анастасия Павлова\par}
\vspace{0.35cm}
{\large \today\par}

\vspace{1.3cm}
{\normalsize Лёвин Артём Александрович эксклюзивно для Анастасии Павловой\par}

\vfill

\begin{tikzpicture}[remember picture,overlay]
  \draw[accent!55,line width=1.25pt]
    ($(current page.south west)+(18mm,17mm)$)
      .. controls
    ($(current page.south west)+(62mm,31mm)$) and
    ($(current page.south east)+(-64mm,8mm)$)
      ..
    ($(current page.south east)+(-18mm,20mm)$);
\end{tikzpicture}

\end{titlepage}

% ================================================================
% ОСНОВНОЙ ЧЕК-ЛИСТ
% ================================================================
\topic{1. Главная идея}

\keyidea{Графики могут быть разными, а общий принцип один: координаты подходящих точек позволяют восстановить неизвестные коэффициенты функции.}

В этом пособии \textbf{узловой точкой} будем называть хорошо считываемую точку графика с удобными, обычно целыми координатами. Если точка $(x_0;y_0)$ принадлежит графику функции $y=f(x)$, то обязательно
\[
y_0=f(x_0).
\]
Именно это равенство превращает координаты точки в уравнение для неизвестных коэффициентов.

\topic{2. Формулы, которые важно узнавать}

\begin{center}
\renewcommand{\arraystretch}{1.35}
\begin{tabularx}{\linewidth}{@{}>{\bfseries}l X X@{}}
\toprule
Вид графика & Формула & Что помнить \\
\midrule
Прямая & $y=kx+b$ & Если неизвестны $k$ и $b$, обычно нужны две информативные точки. \\
Парабола & $y=ax^2+bx+c$ & Число нужных точек зависит от числа неизвестных коэффициентов. \\
Корневая функция & $y=k\sqrt{x}$ & $x\ge0$; точка $(0;0)$ не позволяет определить $k$. \\
Дробно-рациональная функция & $y=\dfrac{k}{x}+a$ & $x\ne0$; особенно внимательно следите за знаком минус. \\
\bottomrule
\end{tabularx}
\end{center}

Если часть коэффициентов уже дана в условии, искать требуется только оставшиеся. Практическое правило:
\[
\boxed{\text{число независимых уравнений}=\text{число неизвестных коэффициентов}.}
\]
Уравнения должны быть информативными и независимыми: формального совпадения числа точек недостаточно.

\topic{3. Линейная функция: $f(x)=kx+b$}

Пусть график проходит через точки
\[
A(-1;-3),\qquad B(3;4).
\]
Подставляем координаты в $f(x)=kx+b$:
\[
\begin{cases}
-3=-k+b,\\
4=3k+b.
\end{cases}
\]
Вычитаем первое уравнение из второго:
\[
4k=7,\qquad k=\frac74.
\]
Затем
\[
-3=-\frac74+b,
\qquad
b=-\frac54.
\]
Следовательно,
\[
\boxed{f(x)=\frac74x-\frac54}.
\]
Например,
\[
f(-5)=\frac74\cdot(-5)-\frac54
=-\frac{40}{4}=-10.
\]

\textbf{Самопроверка:} подставьте $x=-1$ и $x=3$ в найденную формулу. Должны получиться соответственно $-3$ и $4$.

\topic{4. Парабола с двумя неизвестными коэффициентами}

Пусть
\[
f(x)=2x^2+bx+c,
\]
а график проходит через
\[
A(-2;-2),\qquad B(1;1).
\]
Коэффициент при $x^2$ уже известен, поэтому требуется найти только $b$ и $c$:
\[
\begin{cases}
-2=2\cdot(-2)^2-2b+c,\\
1=2\cdot1^2+b+c.
\end{cases}
\]
После упрощения:
\[
\begin{cases}
-2b+c=-10,\\
b+c=-1.
\end{cases}
\]
Отсюда
\[
3b=9,\qquad b=3,\qquad c=-4.
\]
Итак,
\[
\boxed{f(x)=2x^2+3x-4}.
\]

Главный вывод: вид графика изменился, а приём остался тем же --- координаты точки подставляются в формулу функции.

\topic{5. Корневая функция: $f(x)=k\sqrt{x}$}

Для квадратного корня обязательно
\[
\boxed{x\ge0}.
\]

Пусть точка $B(4;5)$ принадлежит графику. Тогда
\[
5=k\sqrt{4}=2k,
\qquad
k=\frac52.
\]
Следовательно,
\[
\boxed{f(x)=\frac52\sqrt{x}}.
\]
Например,
\[
f(6{,}76)=\frac52\sqrt{6{,}76}
=\frac52\cdot2{,}6
=6{,}5.
\]

\textbf{Неудачный выбор точки.} Если взять $(0;0)$, получится
\[
0=k\sqrt0=0.
\]
Это верное равенство, однако значение $k$ из него определить невозможно. Нужно выбрать другую узловую точку.

\topic{6. Дробно-рациональная функция и знак минус}

Рассмотрим
\[
f(x)=\frac{k}{x}+a,
\qquad
\boxed{x\ne0}.
\]
Пусть график проходит через точки $(3;2)$ и $(-3;0)$. Получаем
\[
\begin{cases}
2=\dfrac{k}{3}+a,\\[1mm]
0=-\dfrac{k}{3}+a.
\end{cases}
\]
Из второго уравнения
\[
a=\frac{k}{3}.
\]
Подставляем в первое:
\[
2=\frac{k}{3}+\frac{k}{3}
=\frac{2k}{3},
\qquad
k=3,
\qquad
a=1.
\]
Следовательно,
\[
\boxed{f(x)=\frac{3}{x}+1}.
\]

Полезная запись знака при $b\ne0$:
\[
\boxed{\frac{a}{-b}=-\frac{a}{b}=\frac{-a}{b}}.
\]
Минус в знаменателе лучше сразу перенести перед дробью или в числитель: так дальнейшие преобразования легче контролировать.

\clearpage

% ================================================================
% БЛОК-СХЕМА 1
% ================================================================
\thispagestyle{fancy}

\begin{center}
{\Large\bfseries\color{darkaccent}
Алгоритм: восстановление формулы по узловым точкам}
\end{center}

\vspace{0.65cm}

\begin{center}
\begin{tikzpicture}[node distance=6.5mm]

  \node[flowstart] (start)
    {Дана функция с неизвестными коэффициентами};

  \node[flowaction, below=of start] (write)
    {Запишите общий вид функции и выделите неизвестные коэффициенты};

  \node[flowaction, below=of write] (count)
    {Определите, сколько независимых уравнений требуется};

  \node[flowaction, below=of count] (choose)
    {Выберите столько удобных узловых точек, сколько нужно уравнений};

  \node[flowdecision, below=9mm of choose] (useful)
    {Каждая точка даёт информативное уравнение?};

  \node[flowaction, below=9mm of useful] (subst)
    {Для каждой точки $(x_i;y_i)$ запишите равенство $y_i=f(x_i)$};

  \node[flowaction, below=of subst] (solve)
    {Решите полученную систему и найдите коэффициенты};

  \node[flowaction, below=of solve] (verify)
    {Подставьте исходные точки в найденную формулу};

  \node[flowstart, below=of verify] (finish)
    {Формула восстановлена};

  \draw[flowarrow] (start) -- (write);
  \draw[flowarrow] (write) -- (count);
  \draw[flowarrow] (count) -- (choose);
  \draw[flowarrow] (choose) -- (useful);

  \draw[flowarrow] (useful) --
    node[right,font=\small]{да} (subst);

  \draw[flowarrow] (subst) -- (solve);
  \draw[flowarrow] (solve) -- (verify);
  \draw[flowarrow] (verify) -- (finish);

  \draw[flowarrow] (useful.west) -- ++(-1.8,0)
    node[above,font=\small,pos=0.45]{нет}
    |- (choose.west);

\end{tikzpicture}
\end{center}

\clearpage

% ================================================================
% ПЕРЕСЕЧЕНИЯ ГРАФИКОВ
% ================================================================
\topic{7. Абсцисса и ордината}

Для точки $(x;y)$:
\[
\boxed{\text{абсцисса}=x},
\qquad
\boxed{\text{ордината}=y=f(x)}.
\]

Если требуется ордината точки пересечения, сначала удобно найти её абсциссу $x$, затем подставить найденное значение в формулу одной из функций.

\topic{8. Главное условие пересечения графиков}

Если графики $y=f_1(x)$ и $y=f_2(x)$ пересекаются, в общей точке их ординаты одинаковы. Поэтому
\[
\boxed{f_1(x)=f_2(x)}.
\]

После решения этого уравнения каждое допустимое значение $x$ подставляют в любую из двух исходных функций:
\[
y=f_1(x)
\quad\text{или}\quad
y=f_2(x).
\]
В общей точке обе формулы обязаны дать одно и то же значение $y$.

\subsection*{Пример: пересечение двух прямых}

Пусть
\[
f_1(x)=-5x-6,
\qquad
f_2(x)=-x+1.
\]
Приравниваем формулы:
\[
-5x-6=-x+1.
\]
Отсюда
\[
-4x=7,
\qquad
x=-\frac74.
\]
Теперь находим ординату:
\[
y=f_2\!\left(-\frac74\right)
=\frac74+1
=\frac{11}{4}
=2{,}75.
\]
Точка пересечения:
\[
\boxed{\left(-\frac74;\frac{11}{4}\right)}.
\]

\begin{center}
\begin{tikzpicture}
\begin{axis}[
  width=0.72\linewidth,
  height=0.72\linewidth,
  axis lines=middle,
  xmin=-4,
  xmax=3,
  ymin=-2,
  ymax=5,
  samples=200,
  clip=false,
  grid=both,
  major grid style={linegray!45},
  minor grid style={linegray!18},
  xlabel={$x$},
  ylabel={$y$},
  tick label style={font=\small},
  label style={font=\small},
  unit vector ratio=1 1
]
\addplot[
  accent,
  line width=1.15pt,
  domain=-2.2:-0.8
] {-5*x-6};

\addplot[
  darkaccent,
  line width=1.15pt,
  domain=-4:3
] {-x+1};

\addplot[
  only marks,
  mark=*,
  mark size=2.4pt,
  darkaccent
] coordinates {(-1.75,2.75)};

\node[
  anchor=south west,
  font=\small
] at (axis cs:-1.75,2.75)
  {$\left(-\frac74;\frac{11}{4}\right)$};

\node[
  anchor=south west,
  font=\small,
  text=accent
] at (axis cs:-2.18,4.55)
  {$f_1$};

\node[
  anchor=south west,
  font=\small,
  text=darkaccent
] at (axis cs:-3.7,4.4)
  {$f_2$};

\end{axis}
\end{tikzpicture}
\end{center}

\topic{9. Пересечение с корневой функцией}

Пусть
\[
y_1=2\sqrt{x+5},
\qquad
y_2=3.
\]
Сначала учитываем область определения первой функции:
\[
x+5\ge0,
\qquad
\boxed{x\ge-5}.
\]
В точке пересечения
\[
2\sqrt{x+5}=3.
\]
Делим обе части на $2$:
\[
\sqrt{x+5}=\frac32.
\]
Обе части неотрицательны, поэтому возведение в квадрат здесь равносильно:
\[
x+5=\frac94,
\qquad
x=-\frac{11}{4}.
\]
Найденное значение удовлетворяет $x\ge-5$. Ордината известна из второй формулы:
\[
y=3.
\]
Следовательно,
\[
\boxed{\left(-\frac{11}{4};3\right)}.
\]

\topic{10. Типичные ошибки и способы самоконтроля}

\begin{itemize}[leftmargin=7mm,itemsep=0.45em]

  \item
  \textbf{Перепутаны координаты точки.}
  Всегда читайте $(x;y)$ и используйте равенство $y=f(x)$.

  \item
  \textbf{Выбрана неинформативная точка.}
  Если подстановка даёт только тождество вроде $0=0$, возьмите другую точку.

  \item
  \textbf{Потерян знак минус в дроби.}
  Перенесите минус перед всей дробью или в числитель до дальнейших преобразований.

  \item
  \textbf{Пропущено ограничение.}
  Для $\sqrt{g(x)}$ требуется $g(x)\ge0$, знаменатель не может быть равен нулю.

  \item
  \textbf{Найден только $x$ точки пересечения.}
  Если нужна ордината или вся точка, вычислите также $y$.

  \item
  \textbf{После возведения в квадрат не выполнена подстановка.}
  Сверьте найденные значения с исходным равенством и областью определения.

  \item
  \textbf{Ошибка в восстановленной формуле.}
  Подставьте в неё исходные узловые точки: это быстрый способ обнаружить неверный знак или арифметическую ошибку.

\end{itemize}

\clearpage

% ================================================================
% БЛОК-СХЕМА 2
% ================================================================
\thispagestyle{fancy}

\begin{center}
{\Large\bfseries\color{darkaccent}
Алгоритм: поиск точки пересечения графиков}
\end{center}

\vspace{0.65cm}

\begin{center}
\begin{tikzpicture}[node distance=7mm]

  \node[flowstart] (start2)
    {Даны два графика или две функции};

  \node[flowdecision, below=9mm of start2] (known2)
    {Формулы обеих функций известны?};

  \node[flowside, right=18mm of known2] (recover2)
    {Восстановите недостающие формулы по узловым точкам};

  \node[flowaction, below=10mm of known2] (domain2)
    {Запишите область определения и другие ограничения};

  \node[flowaction, below=of domain2] (equal2)
    {Составьте уравнение $f_1(x)=f_2(x)$};

  \node[flowaction, below=of equal2] (solve2)
    {Решите уравнение относительно $x$};

  \node[flowaction, below=of solve2] (check2)
    {Оставьте только значения, допустимые в исходных функциях};

  \node[flowdecision, below=9mm of check2] (valid2)
    {Остались допустимые значения $x$?};

  \node[flowaction, below=9mm of valid2] (findy2)
    {Для каждого допустимого $x$ найдите $y$ по любой исходной функции};

  \node[flowstart, below=of findy2] (finish2)
    {Запишите точку или точки пересечения};

  \node[flowside, right=18mm of valid2] (none2)
    {Точек пересечения нет};

  \draw[flowarrow] (start2) -- (known2);

  \draw[flowarrow] (known2) --
    node[right,font=\small]{да} (domain2);

  \draw[flowarrow] (known2.east) --
    node[above,font=\small]{нет} (recover2.west);

  \draw[flowarrow] (recover2.south) |- (domain2.east);

  \draw[flowarrow] (domain2) -- (equal2);
  \draw[flowarrow] (equal2) -- (solve2);
  \draw[flowarrow] (solve2) -- (check2);
  \draw[flowarrow] (check2) -- (valid2);

  \draw[flowarrow] (valid2) --
    node[right,font=\small]{да} (findy2);

  \draw[flowarrow] (findy2) -- (finish2);

  \draw[flowarrow] (valid2.east) --
    node[above,font=\small]{нет} (none2.west);

\end{tikzpicture}
\end{center}

\clearpage

% ================================================================
% ТРЕНИРОВОЧНЫЙ БЛОК
% ================================================================
\topic{Тренировочный блок --- Путь А}

\textbf{Пять заданий.}
Оформляйте ключевые подстановки и системы уравнений.
В заданиях с корнями и дробями обязательно указывайте область определения.

\begin{enumerate}[
  leftmargin=8mm,
  label=\textbf{\arabic*.},
  itemsep=1.2em
]

  \item
  Линейная функция
  \[
  f(x)=kx+b
  \]
  проходит через точки $A(-2;5)$ и $B(2;-3)$.
  Найдите формулу функции и вычислите $f(-1)$.

  \item
  Функция
  \[
  f(x)=2x^2+bx+c
  \]
  проходит через точки $A(-1;-3)$ и $B(2;6)$.
  Найдите $b$, $c$ и значение $f(3)$.

  \item
  График функции
  \[
  f(x)=k\sqrt{x}
  \]
  проходит через точку $B(4;6)$.
  Найдите $k$, укажите область определения функции и найдите абсциссу точки пересечения этого графика с прямой $y=9$.

  \item
  Функция
  \[
  f(x)=\frac{k}{x}+a
  \]
  проходит через точки $A(2;4)$ и $B(-2;0)$.
  Найдите $k$, $a$, область определения и значение $f(4)$.

  \item
  Первая прямая проходит через точки
  \[
  A(-1;4),\qquad B(2;-2),
  \]
  вторая --- через точки
  \[
  C(0;-1),\qquad D(3;5).
  \]
  Восстановите формулы обеих прямых и найдите координаты их точки пересечения.

\end{enumerate}

\clearpage

% ================================================================
% ОТВЕТЫ
% ================================================================
\topic{Ответы}

\renewcommand{\arraystretch}{1.55}

\begin{center}
\begin{tabularx}{\linewidth}{@{}c X@{}}
\toprule
\textbf{№} & \textbf{Ответ} \\
\midrule

1 &
$f(x)=-2x+1$, \quad $f(-1)=3$.
\\

2 &
$b=1$, \quad $c=-4$, \quad $f(3)=17$.
\\

3 &
$k=3$, \quad $x\ge0$, \quad абсцисса точки пересечения: $x=9$.
\\

4 &
$k=4$, \quad $a=2$, \quad $x\ne0$, \quad $f(4)=3$.
\\

5 &
$y_1=-2x+2$, \quad
$y_2=2x-1$, \quad
$\left(\dfrac34;\dfrac12\right)$.
\\

\bottomrule
\end{tabularx}
\end{center}

\vfill

\begin{center}
{\small\color{linegray}
Сначала восстановите формулу, затем выполняйте требуемое действие.}
\end{center}

\end{document}